% basic nastaveni
\documentclass[12pt]{article}
\setlength{\parindent}{0pt}
\setlength{\parskip}{1em}
\usepackage[utf8]{inputenc}
\usepackage[czech]{babel}
\usepackage[T1]{fontenc}
\usepackage{graphics}
\usepackage{caption}
\usepackage{hyperref}
\usepackage{dsfont}
\usepackage{color}
\usepackage{float}
\usepackage{enumitem}
\usepackage{graphicx}

% matematicke package
\usepackage{tikz-cd}
\usetikzlibrary{babel}
\usepackage{amsmath}
\usepackage{amssymb}
\usepackage{amsfonts}
\usepackage{amssymb}
\usepackage{geometry}
\geometry{margin=2.5cm}
\usepackage{pgfplots}
\pgfplotsset{compat=1.18}
\usepackage{mathrsfs}
\usetikzlibrary{angles,quotes}

\title{Maturitní otázka číslo 13: Základy planimetrie a konstrukční úlohy}
\author{}
\date{}

\begin{document}
\maketitle
\section{Základní pojmy}
\begin{itemize}
    \item \textbf{Bod, přímka a rovina} jsou nedefinované pojmy(chápeme je intuitivně) a z nich odvozujeme vše ostatní
    \item \textbf{Úsečka} je část přímky ohraničená 2 body
    \item \textbf{Polopřímka} je část přímky, která je ohraničená jen jedním bodem, je tedy nekonečná
    \item \textbf{Polorovina} je část roviny, která je ohraničená jednou přímkou
    \item \textbf{Úhel} je část roviny vymezená dvěma polopřímkami se společným počátkem.
    \item \textbf{Vedlejší úhly} jsou úhly, která leží na stejné přímce a vzájemně se doplňují do $180^\circ$, v obrázku jsou to například úhly $\alpha$ a $\omega$
    \item \textbf{Souhlasné úhly} leží n stejné straně příčky i rovnoběžek, v obrazku jsou to úhly $\alpha , \alpha'$ a $\beta, \beta'$
    \item \textbf{Střídavé úhly} leží na opačných stranách a jsou shodné, v obrázku to jsou úhly $\beta$ a $\gamma'$, nebo úhly $\alpha'$ a $\delta $.
    \item \textbf{Vrcholové úhly} jsou shodné a v obrázku je to například dvojice $\alpha$ a $\delta$.
\end{itemize}


\begin{tikzpicture}[>=stealth, scale=1]
    % Definice přímek
    \draw[thick] (0,2) -- (4,2) node[right] {$p$};
    \draw[thick] (0,0) -- (4,0) node[right] {$q$};
    \draw[thick, blue] (1,-1) -- (3,3) node[above] {$s$};

    % Definice bodů (průsečíků)
    \coordinate (A) at (2.5,2);
    \coordinate (B) at (1.5,0);
    \coordinate (P1) at (3.5,2);
    \coordinate (P2) at (4,5);
    \coordinate (Q1) at (2.5,0);
    \coordinate (Q2) at (1,-1);
    \coordinate (Q3) at (3,0);
    \coordinate (Q4) at (0, 0);
    \coordinate (P3) at (0, 2);
    \coordinate (P4) at (4,2);

    % Vykreslení úhlů
    % Souhlasné úhly (např. vpravo nahoře)
    \draw pic [draw, fill=red!30, angle radius=0.6cm, "$\alpha$"] {angle = P1--A--P2};
    \draw pic [draw, fill=red!30, angle radius=0.6cm, "$\alpha'$"] {angle = Q1--B--A};

    % Střídavé úhly (např. vnitřní)
    \draw pic [draw, fill=green!30, angle radius=0.5cm, "$\beta$"] {angle = B--A--P1};
    \draw pic [draw, fill=green!30, angle radius=0.5cm, "$\beta'$"] {angle = Q2--B--Q3};
    %gamma uhly
    \draw pic [draw, fill=blue!30, angle radius = 0.55cm, "$\gamma$"] {angle = P2--A--P3};
    \draw pic [draw, fill=blue!30, angle radius = 0.55cm, "$\gamma'$"] {angle = P2--B--Q4};
    %delta uhly
    \draw pic [draw, fill=yellow!30, angle radius = 0.65cm, "$\delta$"] {angle = P3--A--B};
    \draw pic [draw, fill=yellow!30, angle radius = 0.65cm, "$\delta'$"] {angle = Q4--B--Q2};

\end{tikzpicture}

\section{Trojúhelník}
Věty o shodnosti: \textbf{sss, sus, usu, Ssu}, u věty Ssu jsou dva trojúhelníky shodné, pokud se shodují ve dvou stranách a úhlu, který leží proti větší z těchto stran
Významné body a přímky:
\begin{itemize}
    \item \textbf{Těžiště} je průsečík těžnic, dělí těžnici v poměru 2:1
    \item \textbf{Kuržnice opsaná} protíná všechny tři vrcholy trojúhelníku a střed je průsečík os stran
    \item \textbf{Kružnice vepsaná} se dotýká stran trojúhelníku a její střed leží v průsečíku os úhlů
    \item \textbf{Výška} vede z vrcholu trojúhelníku na protější stranu, se kterou je kolmá.
    \item \textbf{Thaletova kružnice} nad průměrem $AB$ je množina všech bodů $C$ v rovině, z nichž je úsečka $AB$ vidět pod pravým úhlem s výjimkou bodů $AB$. Středem této kružnice je střed úsečky $AB$. Je to speciální případ věty o středovém a obvodovém úhlu.
    \begin{tikzpicture}[>=stealth, scale=1]
    % Definice přímek
    \draw[thick] (0,0) -- (0.8,1.6) node[right] {$$$$};
    \draw[thick] (0.8,1.6) -- (4,0) node[right] {$$};
    \draw[thick] (0,0) -- (4,0) node[above] {$$};
    \draw (2,0) circle (2cm);

    % Definice bodů (průsečíků)
    \coordinate[label = above:$C$] (C1) at (0.8,1.6);
    \coordinate[label = left:$A$] (A1) at (0,0);
    \coordinate[label = right:$B$] (B1) at (4,0);
    %zvýraznení bodů
    \filldraw (4,0) circle (2pt);
    \filldraw (0,0) circle (2pt);
    \filldraw (0.8, 1.6) circle (2pt);
    % Vykreslení úhlů
    \draw pic [draw, fill=red!30, angle radius = 0.5cm, "$\gamma$"] {angle = A1--C1--B1};


    % Střídavé úhly (např. vnitřní)

    %gamma uhly

    %delta uhly


\end{tikzpicture}

    \textbf{Důkaz thaletovy kružnice:}
    Upravme si původní obrázek takto:\\
    \begin{tikzpicture}[>=stealth, scale=1.5]


    % Definice bodů (průsečíků)
    \coordinate[label = above:$C$] (C1) at (0.8,1.6);
    \coordinate[label = left:$A$] (A1) at (0,0);
    \coordinate[label = right:$B$] (B1) at (4,0);
    \coordinate[label = above:$S$] (S1) at (2,0);



    % Vykreslení úhlů
    \draw pic [draw, fill=red!30, angle radius = 0.7 cm, "$\beta$"] {angle = S1--C1--B1};
    \draw pic [draw, fill=red!30, angle radius = 0.7 cm, "$\beta$"] {angle = C1--B1--A1};
    \draw pic [draw, fill=red!30, angle radius = 0.6 cm, "$\alpha$"] {angle = B1--A1--C1};
    \draw pic [draw, fill=red!30, angle radius = 0.6 cm, "$\alpha$"] {angle = A1--C1--S1};

        %zvýraznění bodů
    \filldraw (4,0) circle (2pt);
    \filldraw (0,0) circle (2pt);
    \filldraw (0.8, 1.6) circle (2pt);
    \filldraw (2,0) circle (2pt);

        % Definice přímek
    \draw[thick] (0,0) -- (0.8,1.6);
    \draw[thick] (0.8,1.6) -- (4,0);
    \draw[thick] (0,0) -- (4,0);
    \draw (2,0) circle (2cm);
    \draw[thick] (0.8,1.6) -- (2,0);

    
\end{tikzpicture}

Platí, že trojúhelníky $SCB$ a $ASB$ jsou rovnoramenné, protože mají vždy dvě strany stejně dlouhé. U trojúhelníku $SCB$ jsou to strany $SC$ a $SB$ protože jejich délka je rovna poloměru kružnice a podobně je to u druhého trojúhleníku. Proto platí, že velikost úhlů $\beta$ je stejná a i velikost úhlů $\alpha$ je stejná.

Dále víme, že součet úhlů v trojúhelníku nám musí vždy dát $180^\circ$. Součet úhlů v trojúhelníku $ABC$ pomocí úhlů $\alpha$ a $\beta$ je tedy:
\[2\alpha + 2\beta = 180^\circ  \qquad \slash \div 2\]
\[\alpha + \beta = 90^\circ\]
A důkaz je hotový. Víme, že součet úhlů $\alpha + \beta$ je rovný $90^\circ$ a víme, že tento součet nám dá zkoumaný vnitřní úhel $ACB$.
\end{itemize}

\end{document}